\clearpage
\SourceLeaf{142}
Es ist der mittlere Werth von $n$ [für $m=5$ und $6$]

\[
(5,1)+2(5,2)+3(5,3)+4(5,4)
\]
\[
= 10x-10xx+5x^3-x^4,
\]
\[
(6,1)+2(6,2)+3(6,3)+4(6,4)+5(6,5)
\]
\[
= 15x-20xx+15x^3-6x^4+x^5
\]

und allgemein

\[
(m,1)+2(m,2)+3(m,3)+\cdots = \frac{(1-x)^m-(1-mx)}{x}.
\]

Bezeichnet man die Zahlencoefficienten zum Beispiel von $(8,5)$ mit
$C^{8.5}$, so ist z.\,B. $3C^{8.5}+C^{8.6}=C^{9.6}$,
[und allgemein

\[
(m-n)C^{m.n}+C^{m.n+1}=C^{m+1.n+1}.
\]

Ferner ist] allgemein

\[
(m-n)x(m,n)+(1-(m-n-1)x)(m,n+1)=(m+1,n+1).
\]

Hieraus folgt leicht, wenn man

\[
(m,1)+2(m,2)+3(m,3)+4(m,4)+\cdots = S_m
\]

setzt, und sich erinnert, dass

\[
(m,0)+(m,1)+(m,2)+(m,3)+\cdots = 1
\]

ist,

\[
S_{m+1} = (1-x)S_m+mx,
\]

und daraus der oben angegebene Werth

\[
\frac{(1-x)^m-(1-mx)}{x}.
\]

\bigskip
\hrule
